\documentclass{amsart}
% For dealing with references we use the comment environment
\usepackage{verbatim}
\newenvironment{reference}{\comment}{\endcomment}
%\newenvironment{reference}{}{}
\newenvironment{slogan}{\comment}{\endcomment}
\newenvironment{history}{\comment}{\endcomment}

% For commutative diagrams we use Xy-pic
\usepackage[all]{xy}

% We use 2cell for 2-commutative diagrams.
\xyoption{2cell}
\UseAllTwocells

% We use multicol for the list of chapters between chapters
\usepackage{multicol}

% This is generally recommended for better output
\usepackage{lmodern}
\usepackage[T1]{fontenc}

% For cross-file-references

% Package for hypertext links:
\usepackage{hyperref}

% For any local file, say "hello.tex" you want to link to please

% Theorem environments.
%
\theoremstyle{plain}
\newtheorem{theorem}[subsection]{Theorem}
\newtheorem{proposition}[subsection]{Proposition}
\newtheorem{lemma}[subsection]{Lemma}

\theoremstyle{definition}
\newtheorem{definition}[subsection]{Definition}
\newtheorem{example}[subsection]{Example}
\newtheorem{exercise}[subsection]{Exercise}
\newtheorem{situation}[subsection]{Situation}

\theoremstyle{remark}
\newtheorem{remark}[subsection]{Remark}
\newtheorem{remarks}[subsection]{Remarks}

\numberwithin{equation}{subsection}

% Macros
%
\def\lim{\mathop{\mathrm{lim}}\nolimits}
\def\colim{\mathop{\mathrm{colim}}\nolimits}
\def\Spec{\mathop{\mathrm{Spec}}}
\def\Hom{\mathop{\mathrm{Hom}}\nolimits}
\def\Ext{\mathop{\mathrm{Ext}}\nolimits}
\def\SheafHom{\mathop{\mathcal{H}\!\mathit{om}}\nolimits}
\def\SheafExt{\mathop{\mathcal{E}\!\mathit{xt}}\nolimits}
\def\Sch{\mathit{Sch}}
\def\Mor{\mathop{\mathrm{Mor}}\nolimits}
\def\Ob{\mathop{\mathrm{Ob}}\nolimits}
\def\Sh{\mathop{\mathit{Sh}}\nolimits}
\def\NL{\mathop{N\!L}\nolimits}
\def\CH{\mathop{\mathrm{CH}}\nolimits}
\def\proetale{{pro\text{-}\acute{e}tale}}
\def\etale{{\acute{e}tale}}
\def\QCoh{\mathit{QCoh}}
\def\Ker{\mathop{\mathrm{Ker}}}
\def\Im{\mathop{\mathrm{Im}}}
\def\Coker{\mathop{\mathrm{Coker}}}
\def\Coim{\mathop{\mathrm{Coim}}}

% Boxtimes
%
\DeclareMathSymbol{\boxtimes}{\mathbin}{AMSa}{"02}

%
% Macros for moduli stacks/spaces
%
\def\QCohstack{\mathcal{QC}\!\mathit{oh}}
\def\Cohstack{\mathcal{C}\!\mathit{oh}}
\def\Spacesstack{\mathcal{S}\!\mathit{paces}}
\def\Quotfunctor{\mathrm{Quot}}
\def\Hilbfunctor{\mathrm{Hilb}}
\def\Curvesstack{\mathcal{C}\!\mathit{urves}}
\def\Polarizedstack{\mathcal{P}\!\mathit{olarized}}
\def\Complexesstack{\mathcal{C}\!\mathit{omplexes}}
% \Pic is the operator that assigns to X its picard group, usage \Pic(X)
% \Picardstack_{X/B} denotes the Picard stack of X over B
% \Picardfunctor_{X/B} denotes the Picard functor of X over B
\def\Pic{\mathop{\mathrm{Pic}}\nolimits}
\def\Picardstack{\mathcal{P}\!\mathit{ic}}
\def\Picardfunctor{\mathrm{Pic}}
\def\Deformationcategory{\mathcal{D}\!\mathit{ef}}

\setlength{\emergencystretch}{3em}
\begin{document}
\title[Stacks Project, Tag 0F3B]{563 --- Weightings of locally quasi-finite morphisms survive arbitrary base change}
\maketitle
\noindent Copyright (C) 2005--2025 Johan de Jong. Stacks Project authors. Source: \href{https://stacks.math.columbia.edu/tag/0F3B}{Tag 0F3B}.

\noindent Editorial difficulty note (not author text). Difficulty H1: A test finite neighbourhood after base change need not come directly from a test neighbourhood on the original target. The proof isolates purely inseparable single-point fibres étale locally and compares the two neighbourhoods on a common étale refinement, eliminating both boundary differences by their closed images..

\noindent Editorial provenance note (not author text). History: excerpt from revision \texttt{a0444\allowbreak 6e57e\allowbreak c1fbc\allowbreak 252a8\allowbreak 71afc\allowbreak ec775\allowbreak 2fb28\allowbreak 07b14}, more-morphisms.tex, lines 22172--22236. Reference presentation was adapted; original-author context and core support blocks were retained or added. The mathematical wording is unchanged. Licensed under GNU FDL 1.2 or later; no invariant sections or cover texts. See ../licenses/COPYING. Context blocks reproduce author definitions and supporting results; they are not additional counted problems.

\begin{definition}
\label{definition-weighting}
Let $f : X \to Y$ be a locally quasi-finite morphism. A
{\it weighting} or a {\it pond\'eration} of $f$ is a map
$w : X \to \mathbf{Z}$ such that for any diagram
$$
\xymatrix{
X \ar[d]_f & U \ar[l]^h \ar[d]^\pi \\
Y & V \ar[l]_g
}
$$
where $V \to Y$ is \'etale, $U \subset X_V$ is open, and $U \to V$ finite,
the function $\int_\pi (w \circ h)$ is locally constant.
\end{definition}

\section{Weightings}
\label{section-weightings}

\noindent
The material in this section is taken from \href{https://stacks.math.columbia.edu/bibliography}{Bibliography: SGA4 (Exposee XVII, 6.2.4)}.

\medskip\noindent
Let $\pi : U \to V$ be a locally quasi-finite morphism of schemes with
finite fibres. Given a function $w : U \to \mathbf{Z}$ we define a function
$$
\textstyle{\int}_\pi w : V \longrightarrow \mathbf{Z},\quad
v \longmapsto
\sum\nolimits_{u \in U,\ \pi(u) = v} w(u) [\kappa(u) : \kappa(v)]_s
$$
Note that the field extensions are finite
(Morphisms, Lemma \href{https://stacks.math.columbia.edu/tag/01TG}{lemma residue field quasi finite (Tag 01TG)}),
$[\kappa' : \kappa]_s$ is the separable degree
(Fields, Definition \ref{fields-definition-insep-degree}),
and the sum is finite as the fibres of $\pi$ are assumed finite.
Another way to compute the value of $\int_\pi w$ at a point $v \in V$
is as follows.
Choose an algebraically closed field $k$ and a morphism
$\overline{v} : \Spec(k) \to V$ whose image is $v$. Then we have
$$
(\textstyle{\int}_\pi w)(v) =
\sum\nolimits_{\overline{u} \in U_{\overline{v}}} w(\overline{u})
$$
where of course $w(\overline{u})$ denotes the value of $w$ at the
image $u$ of the point $\overline{u}$ under the morphism
$U_{\overline{v}} \to U$.
Note that we may view $\overline{u} \in U_{\overline{v}}$ as morphisms
$\overline{u} : \Spec(k) \to U$ such that
$\pi \circ \overline{u} = \overline{v}$. Namely, since
$U \to V$ is locally quasi-finite with finite fibres,
the scheme $U_{\overline{v}}$ is the spectrum
of a finite dimension algebra over $k$ and all of whose prime ideals
are maximal ideals with residue field $k$. To see that the equality
holds, note that the number of morphisms $\overline{u}$ lying over
a given $u$ is equal to $[\kappa(u) : \kappa(v)]_s$ by
Fields, Lemma \href{https://stacks.math.columbia.edu/tag/09HJ}{lemma separable degree (Tag 09HJ)}.



\begin{definition}
\label{fields-definition-insep-degree}
Let $E/F$ be an algebraic field extension. Let $E_{sep}$ be the subextension
found in Lemma \href{https://stacks.math.columbia.edu/tag/030K}{lemma separable first (Tag 030K)}.
\begin{enumerate}
\item The integer $[E_{sep} : F]$ is called the {\it separable
degree} of the extension. Notation $[E : F]_s$.
\item The integer $[E : E_{sep}]$ is called the {\it inseparable
degree}, or the {\it degree of inseparability} of the extension.
Notation $[E : F]_i$.
\end{enumerate}
\end{definition}

\begin{lemma}
\label{lemma-weighting-base-change}
Let $f : X \to Y$ be a locally quasi-finite morphism.
Let $w : X \to \mathbf{Z}$ be a weighting. Let $f' : X' \to Y'$
be the base change of $f$ by a morphism $Y' \to Y$. Then the
composition $w' : X' \to \mathbf{Z}$ of $w$ and the projection $X' \to X$
is a weighting of $f'$.
\end{lemma}

\begin{proof}
Consider a diagram
$$
\xymatrix{
X' \ar[d]_{f'} & U' \ar[l]^{h'} \ar[d]^{\pi'} \\
Y' & V' \ar[l]_{g'}
}
$$
as in Definition \ref{definition-weighting} for the morphism $f'$.
For any $v' \in V'$ we have to show that $\int_{\pi'} (w' \circ h')$ is
constant in an open neighbourhood of $v'$. By
Lemma \href{https://stacks.math.columbia.edu/tag/0F39}{lemma weighting check after etale base change (Tag 0F39)}
(and the fact that \'etale morphisms are open)
we may replace $V'$ by any \'etale neighbourhood of $v'$.
After replacing $V'$ by an \'etale neighbourhood of $v'$
we may assume that $U' = U'_1 \amalg \ldots \amalg U'_n$
where each $U'_i$ has a unique point $u'_i$ lying over $v'$
such that $\kappa(u'_i)/\kappa(v')$ is purely inseparable,
see Lemma \href{https://stacks.math.columbia.edu/tag/02LO}{lemma etale splits off quasi finite part technical variant (Tag 02LO)}.
Clearly, it suffices to prove that $\int_{U'_i \to V'} w'|_{U'_i}$
is constant in a neighbourhood of $v'$.
This reduces us to the case discussed in the next paragraph.

\medskip\noindent
We have $v' \in V'$ and there is a unique point $u'$ of $U'$
lying over $v'$ with $\kappa(u')/\kappa(v')$ purely inseparable.
Denote $x \in X$ and $y \in Y$ the image of $u'$ and $v'$.
We can find an \'etale neighbourhood $(V, v) \to (Y, y)$
and an open $U \subset X_V$ such that $\pi : U \to V$ is finite
and such that there is a unique point $u \in U$ lying over $v$
which maps to $x \in X$ via the projection $h : U \to X$
such that moreover $\kappa(u)/\kappa(v)$ is
purely inseparable. This is possible by the lemma used above.
Consider the morphism
$$
U'' = U \times_X U' \longrightarrow V \times_Y V' = V''
$$
Since $u$ and $u'$ both map to $x \in X$ there is a point
$u'' \in U''$ mapping to $(u, u')$. Denote $v'' \in V''$
the image of $u''$. After replacing $V', v'$ by $V'', v''$ 
we may assume that the composition $V' \to Y' \to Y$ factors
through a map of \'etale neighbourhoods $(V', v') \to (V, v)$
such that the induced morphism $X'_{V'} = X_{V'} \to X_V$ sends
$u'$ to $u$. Inside the base change $X'_{V'} = X_{V'}$ we have
two open subschemes, namely $U'$ and the inverse image $U_{V'}$ of
$U \subset X_V$. By construction both contain a unique point lying
over $v'$, namely $u'$ for both of them.
Thus after shrinking $V'$ we may assume these open subsets
are the same; namely, $U' \setminus (U' \cap U_{V'})$ and
$U_{V'} \setminus (U' \cap U_{V'})$ have
a closed image in $V'$ and these images do not contain $v'$.
Thus $U' = U_{V'}$ and we find a cartesian diagram as in
Lemma \href{https://stacks.math.columbia.edu/tag/0F39}{lemma weighting check after etale base change (Tag 0F39)}.
Since $\int_\pi (w \circ h)$ is locally constant
by assumption we conclude.
\end{proof}
\end{document}
